% !TeX program = xelatex
\documentclass{article}
\usepackage{kotex}
\setmainhangulfont{HCR Batang}
\usepackage{setspace}  % 줄간격
\setstretch{1.3}
\usepackage{enumitem}  % 목록 들여쓰기 
\setlist[itemize]{leftmargin=*, align=left}
\setlist[description]{
  font=\normalfont,
  labelwidth=4cm,
  leftmargin=4cm,
  labelsep=0cm,
  align=left
}
\usepackage{titlesec}  % 섹션 스타일 조정
\titleformat{\section}
  {\normalfont\Large\bfseries}
  {}
  {0pt}
  {}

\titlespacing{\section}
  {0pt}
  {1.5em}
  {0.8em}
\newcommand{\descriptionrule}{%
  \item[]\vspace{-0.6em}
  \noindent\hspace*{-\leftmargin}%
  \rule{0.6\linewidth}{0.4pt}
  \vspace{-0.4em}
}              
\usepackage[a4paper, margin=2.5cm]{geometry}
\usepackage{fancyhdr}                 
\pagestyle{fancy}
\fancyhf{}                             
\fancyfoot[L]{\nouppercase{2026-2, 음악과 사회(002)}} 
\fancyfoot[R]{\thepage}                
\renewcommand{\headrulewidth}{0pt}    
\renewcommand{\footrulewidth}{0.4pt} 
\usepackage[hidelinks]{hyperref}
\let\plainhref\href
\renewcommand{\href}[2]{\plainhref{#1}{\underline{#2}}}

\begin{document}
\title{내 인생의 사운드트랙}
\author{2026학년도 2학기 `음악과 사회' 중간 과제}
\date{언어학과 이승헌(2025-17960)}

\setcounter{secnumdepth}{0}
\bigskip

\maketitle
\thispagestyle{fancy}

\section{1. 들어가며}

\indent 현대인에게 음악 없는 삶이란 상상하기 어렵다. 우리는 출근하거나 이동할 때, 혹은 무언가에 집중할 때, 그밖의 수많은 순간을 음악과 함께 한다. 그런 점에서, 살아가며 음악에 관한 고민을 한다는 것은(예컨대, 어떤 노래를 들을지 고민하는 일, 노래를 듣고 깊은 감정을 느끼는 일 등) 중요하고 숭고한 고민이자 소중한 경험이다. 나 역시 일상의 많은 순간을 음악과 함께해 왔다. 어떤 곡은 평범한 일상의 배경음이 되었고, 어떤 곡은 무대에서 느꼈던 긴장과 설렘을 다시 불러냈다. 또 어떤 곡은 야구장에서 수많은 사람과 한목소리로 부르던 기억과 이어져 있다. 이처럼, 우리네 음악 취향에는 우리의 과거와 현재가 녹아 있다. 이 글에서는 내 삶에 깊은 인상을 남긴 음악을 주제별로 되짚어 보고, 각 곡이 나에게 어떤 의미로 남아 있는지 살펴보고자 한다.

\section{2. 일상과 음악}
\label{sec.2}

\indent 평범한 일상 속에서, 음악이 우리에게 반드시 필요한 존재는 아니라고 생각한다. 그럼에도 나는 일상의 많은 순간에 음악을 듣는다. 혼자 이동하거나 공부를 시작하기 전과 같은 사소한 순간에도 나의 이어폰은 자연스럽게 음악을 재생하고 있다. 때로는 그날의 날씨와 기분에 어울리는 곡을 찾아 듣고, 때로는 한 곡을 통해 과거의 특정한 장소나 사람을 떠올린다. 반드시 있어야 하는 존재는 아니지만, 함께 하면 우리의 평범한 하루를 조금 특별하게, 다채롭게 해주는 존재. 많은 사람이 음악과 함께 일상을 보내는 것에는 이러한 이유가 있다고 생각한다. 똑같은 하루도 어떤 음악을 듣느냐에 따라 다르게 기억될 수 있다는 점에서, 일상 속 음악은 시간을 감각하고 기억하는 하나의 방식이기도 하다. 지금처럼 가을의 한복판을 지나는 시기에, 나는 넬(NELL)의 \href{https://youtu.be/0bTCIbyvBBc?si=3X3nqA629AnmnH9n}{〈기억을 걷는 시간〉}을 찾아 듣는다. 이 곡은 내가 넬을 좋아하게 된 계기이자, 대학에 들어와 밴드의 키보디스트로 활동하고 싶다고 마음먹게 한 곡이다. 넬의 대표곡을 넘어 국내 인디씬을 대표하는 곡이기 때문에 많은 아티스트들이 리메이크하였지만, 그 어떤 버전도 김종완의 음색이 담긴 원곡을 이기지는 못한다.\footnote{여러 리메이크 가운데 JTBC 예능 프로그램 『바라던 바다』에서 샤이니 온유와 블랙핑크 로제가 함께 부른 버전을 추천하고 싶다.} ``아직도'', 세 글자로 시작되는 이 곡의 도입부는 한국 가요 역사상 최고의 도입부라고 칭해도 과언이 아닐 것이다.   
\medskip

 \indent 그렇다면, 가장 인상적인 `전주'를 지닌 곡은 무엇일까. 정말 좋은 곡이 많겠지만, 나는 서태지의 \href{https://youtu.be/WH0iiU0cv00?si=09eEQ3rWFntj04CE}{〈소격동〉}을 꼽을 것 같다. 잔잔한 신디사이저에 이어 울려 퍼지는 높은 음역대의 키보드 리프, 이 곡의 전주가 이어폰을 통해 귀에 꽂히면 이유모를 설렘이 나를 찾아온다. 아이유가 부른 버전도 훌륭하지만, 이 곡의 낯설고 몽환적인 분위기는 원곡자인 서태지의 보컬과 더욱 잘 어울린다고 생각한다. 
\medskip

 \indent 어떤 곡은 영화/드라마 등 다른 매체와 함께 소비되기도 한다. 나에게는 오쿠 하나코(Oku Hanako)의 \href{https://youtu.be/nFG3l5zxLdM?si=ev76Dob4Z8KvGms0}{〈Kawaranai Mono〉}가 그렇다.\footnote{한국어로 번역하면 `변하지 않는 것'이라는 뜻이다.} 이 곡은  영화 『시간을 달리는 소녀』의 OST로 사용된 곡으로, 영화의 순수하고 밝은 분위기와 색감을 떠올리게 하는 곡이다. 영화를 본 뒤 노래를 다시 들으며 가사를 곱씹으면, 변하지 않고 싶은 마음과 그와 무관하게 흘러가는 시간을 다루는 영화의 내용이 더욱 깊게 와닿는다. 이 곡을 들을 때면 중학교 시절 교내 밴드부 친구들과 방과 후에 함께 연주하던 때가 떠오른다. 친구들과 음악으로 즐거움을 나누었던 기억이 담긴 곡이기에, 영화의 감동과는 또 다른 방식으로 마음이 간다. 이처럼 어떤 음악은 과거의 한순간을 다시 불러오는 계기가 되기도 한다.

\section{3. 무대와 음악}
\label{sec.3}

\indent 음악을 소비하는 일과 무대에서 직접 음악에 참여하는 일은 정말 다른 경험으로 다가온다. 청자로서 곡을 들을 때에는 미처 의식하지 못했던 부분이 공연자로서 곡을 들을 때에는 훨씬 구체적으로 다가온다. 무대는 그간 노력해 온 시간을 관객에게 펼쳐 보이는 자리이다. 준비 과정에서 찾아온 어려움은 이후 성장할 수 있는 계기가 되며, 성공적으로 마친 무대는 또 다른 무대를 준비하는 힘이 된다. 내게 무대 위의 음악은 완성된 곡을 재현하는 일을 넘어, 끊임없이 배우며 한계를 넓혀 가는 경험이었다. 이 장에서는 직접 연주한 곡 가운데 무대에서의 긴장과 어려움, 성장을 선명하게 떠올리게 하는 곡들을 골랐다. 먼저 드뷔시의 \href{https://youtu.be/97_VJve7UVc?si=0hGs_9ZKxS90MSxx}{〈Clair de Lune〉}은 내가 처음으로 무대에서 연주한 곡이다. 첫 무대였던 만큼 단순히 음을 연주하는 데 그치지 않고, 악보에 적힌 표현과 곡의 흐름을 세세하게 분석하며 준비했다. 당시에는 무대 경험이 많지 않아 한 음도 틀리지 않아야 한다는 생각이 컸지만, 이 곡을 준비하며 연주는 정확함만으로 완성되지 않는다는 점을 조금씩 깨달았다. 달빛을 떠올리게 하는 섬세한 분위기를 어떻게 전달할지 고민한 과정은 이후 다른 곡을 연주할 때에도 하나의 기준으로 남았다.
\medskip

 \indent 나는 대학에 입학하여 서울대학교 언어학과 밴드 `기표와 기의'에서 여러 차례 공연을 올렸다. 어렸을 때 클래식 피아노를 주로 치던 나에게, 밴드 키보디스트라는 역할은 많은 어려움과 교훈을 주었다. 40여개의 곡을 무대에서 공연했는데, 그 중 준비 과정이 가장 어려웠던 곡은 도쿄지헨(東京事変)\footnote{한국에서는 한자를 그대로 읽어 `동경사변'이라고 부르기도 한다.}의 \href{https://youtu.be/4tlUwgtgdZA?si=3tPUqlSChq6u4uQb}{〈Marunouchi Sadistic〉}이다. 재즈풍의 리듬 속에서 악기끼리 합을 맞추고, 곡 안에서 키보드가 차지할 비중을 고민하는 일이 모두 쉽지 않았다. 그러나 어려움이 컸던 만큼 공연을 통해 많은 것을 배웠다. 특히 밴드 연주는 각자의 파트를 정확하게 연주하는 것만으로 충분하지 않으며, 다른 연주자의 소리를 들으면서 하나의 흐름을 만들어 가야 한다는 점을 실감했다. 이 곡은 내가 키보디스트로서 한 걸음 성장하는 계기가 되었다.
\medskip

 \indent 신해철의 \href{https://youtu.be/qRtgglIu5KY?si=a3lN-o80CFhK-phu}{〈재즈 카페〉}는 화려한 키보드 솔로가 많아 언젠가 무대에서 꼭 연주해 보고 싶었던 곡이다. 지난 겨울에 무대에 올릴 수 있었으며, 지금까지 무대에 올린 곡 가운데 가장 공들여 준비한 곡으로 남아있다. 그동안 많은 곡을 공연하며 배운 점들을 최대한 활용하고자 했고, 그러면서도 클래식 피아노를 통해 단련한 스킬을 녹여내고 싶었다. 단순히 어려운 부분을 반복하는 데 그치지 않고, 키보드 파트가 자연스럽게 어우러지기 바랐다. 키보드를 연주하면서 곡 초반의 랩도 직접 맡았다. 무대에서 마이크를 사용한 것이 처음이라 특히 긴장했고, 연주와 보컬을 동시에 준비해야 한다는 부담도 컸다. 그러나 그만큼 공연을 마친 뒤의 성취감도 컸다. 이 곡은 무대에서 나의 한계를 조금 넘어선 경험과 연결되어 있어 유독 기억에 남는다.

\section{4. 야구와 음악}
\label{sec.4}

\indent 야구장에서 음악은 경기의 배경에 머물지 않는다. 응원가는 관중의 목소리를 하나로 모으고, 선수의 등장곡은 짧은 순간에도 그 선수만의 이미지를 선명하게 만든다. 경기 뒤에 방송되는 음악 역시 승리의 흥분이나 패배의 아쉬움을 정리하며 야구팬의 하루를 마무리한다. 10년이 넘는 시간 동안 야구를 사랑하고 있는 나에게, 음악은 야구에 관한 추억을 떠올리게 하는 중요한 요소이다. 어떤 곡을 들으면 특정 선수의 투구 장면이 떠오르고, 어떤 곡은 이제 사라진 야구장의 풍경을 되살린다. 이 장에서는 응원가와 등장곡, 방송 음악을 통해 야구와 함께 쌓아 온 기억을 돌아보고자 한다. 2015년의 어느 여름날, 대구의 숨 막히는 더위를 뚫고 대구시민운동장 야구장에서 푸른 유니폼의 선수들이 각자의 이름이 적힌 유니폼을 입고 그라운드를 누비고 있었다. 그때, 홈팀 삼성 라이온즈의 응원석에서 경쾌한 전자음 리프가 울려 흘러 나온다. 2001년부터 삼성 라이온즈의 대표 응원가로 사용된 정수라의 곡, \href{https://youtu.be/gJ7h7DPhz_E?si=h7m65UhjbibwVouM}{〈환희〉}다. 매경기의 시작을 알리는 이 응원가는, 나를 항상 설레게 한다. 관중석에서 울려 퍼지는 〈환희〉의 전주가 시작되면, 그날의 경기가 본격적으로 시작된다는 기대감이 야구장 전체를 감싼다. 내 인생에서 절대 빼놓을 수 없는 부분이 야구라면, 야구에서 빼놓을 절대 빼놓을 수 없는 음악이 바로 이 곡일 것이다.  
\medskip

 \indent 넥스트(N.EX.T)의 \href{https://youtu.be/8GKE823m-hw?si=nczjFQG77hKrVEq_}{〈Lazenca, Save Us〉}. 기성세대에게는 TV 애니메이션 『영혼기병 라젠카』의 주제가로, 젊은 세대에게는 가수 하현우의 커버곡으로 널리 알려져 있는 곡이다. 그러나 야구팬들에게 이 노래가 지니는 의미는 사뭇 다르다. 오랫동안 국내 야구를 지켜본 팬들에게 이 곡은 2010년대 리그를 지배했던 삼성 라이온즈의 마무리 투수, 오승환의 등장곡으로 기억된다. 두 번 씩 울리는 베이스드럼, 그 위로 나오는 신해철의 중후한 목소리는 상대팀의 사기를 완전히 앗아 갔다. 그만큼 전성기 오승환의 퍼포먼스가 압도적이었기 때문인데, 〈Lazenca, Save Us〉가 나오면 상대팀팬들은 모두 집에 가곤 했다는 이야기가 있을 정도다. 활발한 응원 문화가 발달한 한국 야구에서 어떤 음악은 특정 선수와 불가분의 관계에 놓이기도 하는데, 이 곡과 오승환의 관계가 꼭 그렇다고 생각한다. 이 곡은 음악이 한 선수의 이미지와 서사를 어떻게 강화하는지를 잘 보여준다.
\medskip

 \indent 2026년, 야구팬들은 잠실야구장의 마지막 순간을 지켜보았다. 40년에 가까운 시간 동안 LG와 두산을 품어 온 잠실야구장은 한국 야구의 메카와 같은 공간이었다. 이곳을 떠올릴 때 빠질 수 없는 음악이 봄여름가을겨울의 \href{https://youtu.be/MhlvBxEW4SQ?si=ok2qn8BEyxltJbPZ&t=18}{〈Bravo, My Life!〉}이다. 2007년부터 두산 베어스의 8회 응원을 책임져 온 이 곡은, 경기 후반에 팬들의 분위기를 끌어올리는 동시에, 잠실 야구장과 함께 했던 수많은 팬들의 기억에 남았다. 당연할 것 같았던 플래시 응원이, 영원할 것 같았던 잠실 야구장이, 이제는 야구팬들의 추억 속 한 장면으로 남게 되었다. 장소는 사라지지만 추억은 사라지지 않을 것이라 믿는다. 그곳에서 울려 퍼진 음악도 마찬가지이지 않을까. 〈Bravo, My Life!〉를 들으며 우리의 잠실 야구장을 다시 한번 추억해 본다. 
\medskip

 \indent 야구 경기가 끝난 뒤에는 그날의 주요 장면을 정리하는 하이라이트 프로그램이 방송된다. 이 프로그램까지 시청해야 비로소 야구팬의 하루가 마무리된다고 해도 과언이 아니다. 마지막으로 소개할 곡은 MBC 스포츠플러스의 하이라이트 프로그램 `베이스볼 투나잇'의 엔딩곡인 마시따밴드의 \href{https://youtu.be/u5Aj7oDDMeQ?si=0zRCUFLrT8UJhXqM}{〈돌멩이〉}이다. 2013년부터 프로그램의 마지막을 장식해 온 이 곡은 작은 돌멩이가 꿈과 희망을 품고 앞으로 나아가는 모습을 노래한다. 이러한 가사 덕분에 경기 안팎의 감동적이거나 극적인 순간을 되짚는 영상의 배경 음악으로 많은 사랑을 받았다. 야구팬의 하루가 마무리되는 순간을 장식한다는 점에서, 〈돌멩이〉는 이 장을 마무리하기에 가장 어울리는 곡인 듯하다.   

\section{5. 삶과 음악}
\label{sec.5}

\indent 삶은 늘 기대한 방향으로만 흘러가지 않기에 때로는 불안과 후회, 조급함을 마주하게 된다. 이럴 때 음악이 문제를 직접 해결해 주지는 않지만, 마음속 감정을 잠시 들여다보고 다시 정리하도록 돕는다. 누군가의 노랫말은 미처 인식하지 못한 생각을 하게 만들고, 때론 혼자 있는 순간에도 외롭지 않은 듯한 기분을 만들어 주기도 한다. 이처럼, 힘든 순간에 찾는 음악은 단순한 기분 전환의 수단이라기보다, 스스로를 다독이고 다시 앞으로 나아가게 하는 작은 버팀목에 가깝다. 이 장에서는 마음이 복잡할 때 나 자신을 응원하기 위해 찾아 듣는 곡들을 소개한다. 먼저 소개할 곡은 국내 인디씬의 라이징스타, 한로로의 \href{https://youtu.be/h0KIWaUEIgQ?si=5VR-VvL-AEm-ooM6}{〈사랑하게 될 거야〉}라는 곡이다. 이 곡은 스스로를(또는 타인을) 미워하기보단 사랑하는 편이 훨씬 낫다는 나의 가치관과 맞닿아 있다. 누구나 자신이 마음에 들지 않거나 스스로를 신뢰하기 어려운 순간을 겪지만, 이 곡은 그런 순간에도 자신을 쉽게 포기하지 말자는 메시지를 건넨다. 내가 나를 사랑하지 못한다고 느낄 때 이 곡을 찾는 이유도 그 때문이다. 내가 나를 사랑해야 하는 이유를 대단히 논리적으로 제시하는 것이 아니라, `당연히 그래야 하는 것'이라는 메시지를 전하는 것이 이 곡의 매력이라고 생각한다. 
\medskip

 \indent 다음 곡은 브릿팝의 아이콘 오아시스(Oasis)의 \href{https://youtu.be/dPPi2D6GK7A?si=uDUmr9USeEactqyv}{〈The Masterplan〉}이다. 이 곡은 주로 학업이나 미래에 대한 막연한 불안과 걱정이 커질 때 자주 듣는 곡이다. 특히 ``We're all part of the masterplan''이라는 가사를 들으면 모든 일을 당장 이해하거나 통제할 수 없더라도, 내가 나를 믿고 열심히 하면 언젠가 좋은 결과를 받아들 수 있을 것이라는 생각이 든다. 대학 입시를 준비하던 시기에 특히 자주 들었으며, 지금도 학업과 진로 문제로 마음이 복잡해질 때마다 자연스럽게 다시 찾게 되는 것 같다. 우린 모두 완벽한 계획의 일부라는 메시지를 전하고 싶다. 
\medskip

 \indent 마지막으로 소개할 곡은 신해철의 \href{https://youtu.be/HRlwPwqC-Y0?si=w3gsKWYJDUgvj0xK}{〈나에게 쓰는 편지〉}다. 대학에 와서 가장 크게 느끼는 것 중 하나는, 나를 둘러 싼 세상이 정말 빠르게 변화한다는 점이다. 누구나 한번쯤은, 세상의 변화로부터 뒤쳐질 수 있겠다는 생각에 불안했던 적이 있을 것이다. 이 곡에는 다음과 같은 가사가 등장한다. ``우리가 찾는 소중함들은 항상 변하지 않아.'' 아무리 세상이 달라지고 시간이 흐른다 해도, 우리가 소중히 여기는 가치까지 모두 사라지는 것은 아니다. 내 주변이 바쁘게 변해도 나는 나를 지킬 수 있다는 믿음, 이 곡은 그러한 믿음에 힘을 실어준다. 과거에도 많이 들었지만, 앞으로 살아가며 더 많이 듣게 될 곡이라고 생각한다. 

\section{6. 나오며}

\indent 내 인생의 사운드트랙을 만들어 보았다. 각기 다른 장르와 시대 속의 곡들이지만, 모두 내 삶의 특정한 장면과 연결되어 있다는 점에서 나에게 남다른 의미를 갖는 곡이다. 어떤 곡은 반복되는 일상의 배경음이 되었고, 어떤 곡은 친구들과 함께 연주했던 시절이나 무대 위의 긴장과 성취를 떠올리게 했다. 또 어떤 곡은 야구장에서 느꼈던 소속감과 열기를 되살리고, 힘든 순간에는 나 자신을 다시 다독이게 한다. 결국 음악은 감상하는 순간에만 머무르지 않고, 그 음악을 들었던 장소와 사람, 감정을 매개하는 수단인 것 같다. 앞으로도 새로운 음악을 만나고 새로운 순간을 살아가겠지만, 이 글에 담은 곡들은 저마다의 방식으로 내 삶을 채워 주는 나만의 사운드트랙으로 오래 남을 것이다.

\newpage
\section{부록. `내 인생의 사운드트랙'}
\bigskip
\begin{itemize}
  \item[\textbf{Track 1.}] 넬 - 〈기억을 걷는 시간〉
  \item[\textbf{Track 2.}] 서태지 - 〈소격동〉
  \item[\textbf{Track 3.}] Oku Hanako - 〈Kawaranai Mono〉
  \vspace{0.5em}
  \item[\textbf{Track 4.}] C. Debussy - 〈Clair de Lune〉
  \item[\textbf{Track 5.}] 東京事変 - 〈Marunouchi Sadistic〉 
  \item[\textbf{Track 6.}] 신해철 - 〈재즈 카페〉
  \vspace{0.5em}
  \item[\textbf{Track 7.}] 정수라 - 〈환희〉
  \item[\textbf{Track 8.}] 넥스트(N.EX.T) - 〈Lazenca, Save Us〉
  \item[\textbf{Track 9.}] 봄여름가을겨울 - 〈Bravo, My Life!〉   
  \item[\textbf{Track 10.}] 마시따밴드 - 〈돌멩이〉
  \vspace{0.5em}
  \item[\textbf{Track 11.}] 한로로 - 〈사랑하게 될 거야〉
  \item[\textbf{Track 12.}] Oasis - 〈The Masterplan〉
  \item[\textbf{Track 13.}] 신해철 - 〈나에게 쓰는 편지〉
\end{itemize}
\end{document}
